\begin{textsample}{NEG Dim 6 – ba – Score 1.00 – ba/ba\_inf\_e\_ef\_8.txt}  \label{ex:f6_neg_213}
O Documento Curricular Referencial da Bahia para a Educação Infantil e Ensino Fundamental contempla e articula os conhecimentos científicos às temáticas da contemporaneidade, por meio dos Temas Integradores, em escala local, regional e global, em uma perspectiva de promover o desenvolvimento de cidadãos autônomos, responsáveis, engajados e imbuídos na formação de uma sociedade mais justa, sustentável, equânime, igualitária, inclusiva e laica.
A seguir são apresentados os Temas Integradores que precisarão ser considerados de forma transversal nos currículos escolares da Educação Básica, em todas as etapas e modalidades do Estado da Bahia. | Temas Integradores | Normativos | Finalidades | |---|---|---| | Educação para a Diversidade — Educação para as Relações de Gênero e Sexualidade | Lei Federal nº 11.340/2006 | Cria mecanismos para coibir a violência doméstica e familiar contra a mulher, nos termos do § 8º do art. 226 da Constituição Federal, da Convenção sobre a Eliminação de Todas as Formas de Discriminação contra as Mulheres e da Convenção Interamericana para Prevenir, Punir e Erradicar a Violência contra a Mulher ; dispõe sobre a criação dos Juizados de Violência Doméstica e Familiar contra a Mulher ; altera o Código de Processo Penal, o Código Penal e a Lei de Execução Penal e dá outras providências. | | Educação para a Diversidade — Educação para as Relações de Gênero e Sexualidade | Lei Federal nº 2.848/40, § 7º ao art. 121 do Código Penal | Estabelece o aumento da pena do feminicídio. | | Educação para a Diversidade — Educação para as Relações de Gênero e Sexualidade | Resolução nº 120/2013 CEE | Dispõe sobre a inclusão do nome social dos/das estudantes travestis, transexuais e outros no tratamento nos registros escolares e acadêmicos nas Instituições de Ensino que integram o Sistema de Ensino do Estado da Bahia e dá outras providências. | | Educação para a Diversidade — Educação para as Relações de Gênero e Sexualidade | Plano Nacional de Políticas para as Mulheres ( 2013-2015 ) | Objetiva o fortalecimento e a institucionalização da Política Nacional para as Mulheres. | | Educação para a Diversidade — Educação para as Relações de Gênero e Sexualidade | Plano Estadual de Políticas para as Mulheres ( 2013-2015 ) | Objetiva o fortalecimento e a institucionalização da Política Estadual para as Mulheres. | | Temas Integradores | Normativos | Finalidades | |---|---|---| | Educação para a Diversidade — Educação das Relações Étnico-Raciais | Lei nº 10.639, de 9 de janeiro de 2003 | Inclui no currículo oficial da Rede de Ensino a obrigatoriedade da temática História e Cultura Afro-Brasileira e dá outras providências. | | Educação para a Diversidade — Educação das Relações Étnico-Raciais | Resolução nº 1/2004 CNE/CEB | Estabelece as Diretrizes Curriculares Nacionais para a Educação das Relações Étnico-Raciais e para o Ensino de História e Cultura Afro-Brasileira e Africana. | | Educação para a Diversidade — Educação das Relações Étnico-Raciais | Lei Federal nº 11.645/2008 | Estabelece as diretrizes e bases da educação nacional, para incluir no currículo oficial da rede de ensino a obrigatoriedade da temática História e Cultura Afro-Brasileira e Indígena. | | Educação para a Diversidade — Educação das Relações Étnico-Raciais | Lei Federal nº 12.288/2010 | Institui o Estatuto da Igualdade Racial ; altera as Leis nº 7.716, de 5 de janeiro de 1989 ; 9.029, de 13 de abril de 1995 ; 7.347, de 24 de julho de 1985, e 10.778, de 24 de novembro de 2003. | | Educação para a Diversidade — Educação das Relações Étnico-Raciais | Lei Estadual nº 13.182/2014 | Institui o Estatuto da Igualdade Racial e de Combate à Intolerância Religiosa do Estado da Bahia. | | Educação em Direitos Humanos | Decreto Presidencial nº 7.037, de 21 de dezembro de 2009 | Aprova o Programa Nacional de Direitos Humanos ( PNDH-3 ) e dá outras providências. | | Educação em Direitos Humanos | Plano Nacional de Educação em Direitos Humanos/2007 | Difunde a cultura de Direitos Humanos no país. | | Educação em Direitos Humanos | Decreto Governamental nº 12.019/2010 | Aprova o Plano Estadual de Direitos Humanos da Bahia ( PEDH ) e dá outras providências. | | Educação em Direitos Humanos | Plano Estadual de Educação em Direitos Humanos/2009 | Expressa o compromisso do Governo do Estado da Bahia com a promoção da cidadania e dos Direitos Humanos. | | Educação em Direitos Humanos | Parecer CEE/CEB nº 8/2012 | Estabelece Diretrizes Nacionais para a Educação em Direitos Humanos. | | Educação Ambiental | Lei Federal nº 9.795/1999 | Institui a Política Nacional de Educação Ambiental e dá outras providências. | | Educação Ambiental | Resolução nº 2/2012 CNE/CP | Estabelece as Diretrizes Curriculares Nacionais para a Educação Ambiental. | | Educação Ambiental | Resolução nº 11/2017 CEE | Dispõe sobre a Educação Ambiental no Sistema Estadual de Ensino da Bahia. | | Educação Ambiental | Lei Estadual nº 12.056/2011 | Institui a Política de Educação Ambiental do Estado da Bahia e dá outras providências. | | Educação Ambiental | Decreto nº 19.083, de 06 de junho de 2019 | Regulamenta a Lei nº 12.056, de 07 de janeiro de 2011, que institui a Política de Educação Ambiental do Estado da Bahia e dá outras providências. | | Temas Integradores | Normativos | Finalidades | |---|---|---| | Saúde na Escola | Lei Federal nº 11.947/2009 | Dispõe sobre o atendimento da alimentação escolar e do Programa Dinheiro Direto na Escola aos alunos da Educação Básica ; altera as Leis nº 10.880, de 9 de junho de 2004, 11.273, de 6 de fevereiro de 2006, 11.507, de 20 de julho de 2007 e revoga dispositivos da Medida Provisória nº 2.178-36, de 24 de agosto de 2001, e a Lei nº 8.913, de 12 de julho de 1994, e dá outras providências. | | Saúde na Escola | Decreto Presidencial nº 6.286/2007 | Institui o Programa Saúde na Escola. | | Saúde na Escola | Portaria Conjunta SEPLAN/SESAB/SEC nº 001/2014 | Institucionaliza as ações transversais e esforços intersetoriais para implantação do Programa de Ação Estadual de Prevenção da gravidez e assistência ao parto na adolescência. | | Saúde na Escola | Portaria nº 2728/2016 | Institui a Promoção da Saúde e Prevenção de Doenças e Agravos no contexto escolar, com ênfase no combate ao mosquito Aedes aegypti. | | Saúde na Escola | Portaria Conjunta SESAB/SEC nº 01/2018 | Dispõe sobre a obrigatoriedade da apresentação da carteira/cartão de vacinação em creches e escolas, em todo o território do Estado da Bahia. | | Educação Fiscal | Portaria Interministerial nº 413/2002 MF/MEC | Implementa o Programa Nacional de Educação Fiscal. | | Educação Fiscal | Decreto Estadual nº 15.737/2014 | Institui a Educação Fiscal na Bahia. | | Educação para o Trânsito | Lei Federal nº 9.503/1997 | Institui o Código de Trânsito Brasileiro. |

% matched lemmas: 
\end{textsample}
